% Transversal construction and the verifier's three kinds of checks.
\begin{frame}{Distillation over the Trap Code}
  \centering
  \setbeamercolor{itemize item}{fg=CVQCTeal}
  \begin{tikzpicture}
    \path[use as bounding box] (-6.6,-3.15) rectangle (6.6,3.05);

    \node[anchor=north west,inner sep=0pt,text=CVQCInk] at (-6.35,2.45) {%
      \begin{minipage}{6.4cm}
        \raggedright
        \begin{itemize}
          \item Compatible with $\mathsf{TEnc}$.
        \end{itemize}
      \end{minipage}%
    };

    \node[anchor=north east,inner sep=0pt,text=CVQCInk] at (6.35,2.65) {%
      \resizebox{5.6cm}{!}{%
        \let\tiny\scriptstyle
        \hspace*{1.7cm}% Include both ciphertext labels in the visual bounds.
        $\Qcircuit @C=1.4em @R=1.2em {
          \lstick{\mathsf{TEnc}(\ket{0})} & \push{\makebox[0pt]{$/$}}\qw & \ctrl{1} & \targ & \qw
            & \rstick{\mathsf{TEnc}(\ket{0})}\qw \\
          \lstick{\mathsf{TEnc}(\ket{0})} & \push{\makebox[0pt]{$/$}}\qw & \targ & \qw & \meter & \cw \\
          \lstick{\mathsf{TEnc}(\ket{0})} & \push{\makebox[0pt]{$/$}}\qw & \ctrl{1} & \ctrl{-2} & \meterB{X} & \cw \\
          \lstick{\mathsf{TEnc}(\ket{0})} & \push{\makebox[0pt]{$/$}}\qw & \targ & \qw & \meter & \cw
        }$\hspace*{1.7cm}%
      }%
    };

    \uncover<2->{
      \node[anchor=north west,fill=CVQCFog,rounded corners=3pt,
          inner xsep=0.3cm,inner ysep=0.25cm,text=CVQCInk] at (-5.8,-0.65) {%
        \begin{minipage}{11cm}
          \raggedright
          \textbf{Reject if:}
          \par\vspace{0.35em}
          \begin{itemize}
            \setlength{\itemsep}{0.15em}
            \item<2-> Any trap check fails.
            \item<3-> Any QEC syndrome is nonzero.
            \item<4-> Any decoded measurement outcome is unexpected.
          \end{itemize}
        \end{minipage}%
      };
    }
  \end{tikzpicture}

  \note<1>{
    \footnotesize
    The circuit uses only CNOTs and Pauli measurements, which have
    transversal implementations over the trap code. Each drawn wire is
    now a logical ciphertext block, indicated by a slash. Apply the C-dist
    CNOTs to every aligned
    physical coordinate; perform physical measurements and let the verifier
    remove the masks, decode, and check the traps and logical outcomes.
    The verifier updates the keys, and all blocks share the hidden
    permutation. The figure shows the honest case and suppresses these
    physical details. Correct encodings give a correct output encoding.
    Next, spell out the verifier's rejection checks. The analysis of
    malformed ciphertexts and malicious execution starts on the next slide.
  }
  \note<2>{
    \footnotesize
    Before these checks, the verifier updates the Pauli keys through the
    Clifford circuit, removes the resulting masks from the reported
    measurement strings, and undoes the hidden permutation. It separates
    the data outcomes from the two trap types. In a Z-basis measurement,
    check the zero traps; in an X-basis measurement, check the plus traps.
    Reject if any of these deterministic trap outcomes is wrong. The
    complementary trap type has random outcomes and is not required to
    have a fixed value. Source: 4-to-1 Distillation Protocol, trap checks.
  }
  \note<3>{
    \footnotesize
    Check the measured QEC data block in the measurement basis. The
    strict codeword-checking decoder rejects a nonzero syndrome rather
    than correcting it. In the paper's protocol, this is the condition
    that the basis-dependent decoder returns bottom. This check is
    distinct from the traps and from the decoded logical outcome.
    Source: CSS code convention and 4-to-1 Distillation Protocol, data
    decoding step.
  }
  \note<4>{
    \footnotesize
    Finally, compare the decoded outcomes with the accepted logical
    measurement set. For zero-state distillation, the Z outcomes on
    copies two and four must be zero. Either X outcome on copy three
    is accepted; there is no deterministic logical check on that copy.
    EPR and Hadamard magic states use the corresponding stabilizer
    parity checks from the paper. If all three kinds of checks pass,
    keep the first ciphertext with its updated key. The output label
    in the figure depicts honest execution; it does not yet assert
    a security guarantee on arbitrary attacked inputs. Next, introduce
    the notation for an input with one plaintext-bad copy, then derive
    its output behavior. Source: 4-to-1 measurement accepted results and
    the 4-to-1 Distillation Protocol.
  }
\end{frame}
